\subsection{오케스트레이션 구조별 자원 효율성 분석}

본 연구에서는 제안한 세 가지 오케스트레이션 구조(Code, LLM, Hybrid)가 문서 생성 과정에서 소모하는 자원 효율성을 정량적으로 분석하였다. 실험 결과, 고정된 그래프 기반의 제어 방식을 사용하는 Code 구조가 자원 소모 측면에서 가장 높은 효율성을 보였다. 표~\ref{tab:perf_main}에 제시된 정형 보고서 생성 작업('금형-사출-사출기-단위-가동품질-통계-분석-보고서')의 실측 데이터에 따르면, Code 구조의 평균 실행 시간은 158.61초, 총 토큰 소모량은 21,440개로 측정되었다.

이러한 결과는 오케스트레이션의 제어 메커니즘이 단순할수록 오버헤드가 감소한다는 가설을 뒷받침한다. Code 구조는 미리 정의된 워크플로우에 따라 에이전트 간의 상태 전이가 결정되므로, 다음 단계를 결정하기 위한 추가적인 LLM 추론 과정이 필요하지 않다. 반면, 동적 라우팅을 수행하는 LLM 구조의 경우, 매 단계마다 현재 상태를 분석하고 최적의 에이전트를 선택하는 판단 과정에서 추가적인 토큰 소모와 지연 시간이 발생하며, 이는 자원 효율성 저하로 이어진다.

\subsection{생성 품질 및 프로세스 안정성 분석}

자원 효율성과는 대조적으로, 산출물의 품질 점수와 재시도 횟수(retry count)는 오케스트레이션의 유연성과 밀접한 상관관계를 보였다. LLM-as-Judge 방식을 통해 산출된 Code 구조의 품질 점수는 평균 82점으로 나타났으며, 최종 결과물을 얻기까지 평균 3회의 재시도가 수행되었다. 이는 고정된 파이프라인이 제공하는 일관된 구조적 완결성 덕분에 일정 수준 이상의 품질을 안정적으로 확보할 수 있음을 의미한다.

그러나 재시도 횟수와 최종 점수 간의 관계를 분석했을 때, 고정된 경로를 따르는 Code 구조는 특정 단계에서 오류가 발생할 경우 이를 수정하기 위해 동일한 경로를 반복 수행해야 하는 한계가 있다. 반면, LLM 구조는 검수 에이전트의 피드백에 따라 유연하게 경로를 수정하는 동적 라우팅이 가능하여, 복잡한 수정 요구사항이 발생하는 작업에서 더 높은 품질 점수를 획득하는 경향을 보였다. 이러한 자원 소모량과 품질 간의 상충 관계(Trade-off)는 그림~\ref{fig:tradeoff}에서 명확히 드러난다. 그림~\ref{fig:tradeoff}의 X축인 자원 소모량과 Y축인 품질 점수의 분포를 살펴보면, Code 구조는 저자원-중품질 영역에 위치하는 반면, LLM 및 Hybrid 구조는 고자원-고품질 영역으로 이동하는 양상을 띤다.

\subsection{작업 성격에 따른 최적 구조 제안}

실험 결과, 작업의 성격에 따라 최적의 오케스트레이션 구조가 상이함이 입증되었다. '금형-사출 분석 보고서'와 같이 엄격한 논리적 구조와 정형화된 데이터 분석이 핵심인 작업에서는 Code 구조가 가장 효율적인 선택지이다. 해당 작업은 설계-집필-검수-수정의 단계가 명확하며, 경로의 유연성보다는 정의된 절차의 충실한 이행이 중요하기 때문에 158.61초의 빠른 실행 시간과 상대적으로 적은 토큰 소모로도 82점이라는 준수한 품질을 달성할 수 있었다.

반면, 정해진 형식이 없거나 문맥적 유연성이 요구되는 창의적 글쓰기 작업의 경우, 고정된 그래프 방식은 생성물의 단조로움을 초래하거나 복잡한 수정 요구사항을 반영하지 못하는 한계를 보인다. 따라서 작업의 성격이 정형 보고서 형태일 때는 제어 오버헤드가 최소화된 Code 구조를, 높은 창의성과 정교한 반복 수정이 필요한 작업일 때는 자원 소모를 감수하더라도 LLM 구조나 계층적 제어를 수행하는 Hybrid 구조를 채택하는 것이 바람직하다. 결과적으로 오케스트레이션 구조의 선택은 단순히 성능의 우열이 아닌, 가용 자원과 요구되는 품질 수준, 그리고 작업의 정형성 정도에 따른 전략적 선택의 문제로 정의할 수 있다.